\section{The Circuit Solver}
\label{sec:circuit-solver}

\subsection{Modified nodal analysis formulation}

The solver follows the standard MNA formulation~\cite{ho1975mna}. With $n$ non-ground nodes
and $m$ independent voltage sources, the unknown vector is the stacked node-voltage vector
$v \in \mathbb{R}^n$ and branch-current vector $j \in \mathbb{R}^m$ (one auxiliary unknown
per independent voltage source, the standard MNA treatment for an ideal source, which
cannot be described by a node-to-node conductance alone):
\begin{equation}
\begin{bmatrix} G & B \\ B^{\top} & 0 \end{bmatrix}
\begin{bmatrix} v \\ j \end{bmatrix}
=
\begin{bmatrix} i \\ e \end{bmatrix},
\end{equation}
where $G$ is the conductance matrix assembled by stamping every resistive (or
companion-resistive, see below) element, $B$ encodes which two nodes each voltage source is
connected across, $i$ collects any independent current injections, and $e$ is the vector of
source values at the current simulation time. Node~0 is reserved, by MNA convention, as the
ground reference and is excluded from the unknown vector entirely; every other node's
voltage is solved for relative to it. The resulting dense linear system is solved once per
game tick (every $1/20$~s of game time) by Gaussian elimination with partial pivoting.

To guard against a genuinely floating node -- a component with an unconnected lead, which a
player building a circuit incrementally will produce constantly and which must not crash
the simulation -- a small conductance ($10^{-9}$~S) is added to every non-ground node's
diagonal unconditionally. This is large enough to keep the matrix non-singular for an
isolated node, and small enough (roughly nine orders of magnitude below the smallest
resistor preset) to be numerically invisible in any circuit that is actually
connected.

\subsection{Reactive elements}

Capacitors and inductors are stamped using trapezoidal-integration companion models, the
same technique used in general-purpose circuit simulators, rather than backward Euler,
specifically because backward Euler introduces artificial numerical damping that would
visibly flatten an underdamped LC transient a student is trying to observe. For a capacitor
with capacitance $C$ and timestep $h$, the companion model is a conductance
$G_{eq} = 2C/h$ in parallel with a Norton current source
$I_{eq} = -G_{eq} v_C(t - h) - i_C(t-h)$, and correspondingly for an inductor,
$G_{eq} = h/(2L)$; both are re-derived from the trapezoidal rule
$x(t) \approx x(t-h) + \frac{h}{2}\left[\dot{x}(t) + \dot{x}(t-h)\right]$ applied to the
element's defining relation ($i = C\,\mathrm{d}v/\mathrm{d}t$ or
$v = L\,\mathrm{d}i/\mathrm{d}t$, respectively).

\subsection{The memristor}

The memristor implements Strukov et al.'s charge-controlled linear-drift
model~\cite{strukov2008missing}: its resistance is
\begin{equation}
R(q) = R_{on}\left(1 - \frac{q}{q_{max}}\right) + R_{off}\,\frac{q}{q_{max}}, \qquad
\dot{q}(t) = i(t),
\end{equation}
with $q/q_{max}$ clamped to $[0,1]$, defaults of $R_{on}=100\,\Omega$ and
$R_{off}=10\,000\,\Omega$, and a player-adjustable ``switching speed'' preset that scales
$q_{max}$. Because $q$ is the time-integral of the current actually flowing through the
device, its resistance depends on the device's entire history rather than its instantaneous
state -- the defining property that distinguishes a memristor from a resistor -- and,
notably, this state is deliberately preserved across a network topology rebuild
(Section~\ref{sec:limitations} discusses why the other reactive elements are not yet
afforded the same treatment).

\subsection{An explicit, player-placed ground reference}

MNA's ground node is, by construction, always present as a mathematical reference: node~0
is never actually solved for and is defined to be exactly $0$~V. Prior to the addition of a
dedicated ground block, however, nothing in the game world ever corresponded to node~0 --
every network's blocks were assigned freshly allocated node indices starting at~1, and
node~0 remained a purely internal bookkeeping device invisible to the player. This meant a
component's \emph{voltage drop} across its own two leads was always well defined (differences
between node voltages are invariant to which node is chosen as the reference), but there was
no way to ask a more basic question: ``what is the voltage \emph{at} this point in the
circuit?'', since that requires an agreed reference point, exactly as it does on a real
bench before a probe's ground clip is attached somewhere.

The ground block is conductive on all six faces, like wire, and is otherwise electrically
inert; the only special handling it receives is in the topology-construction step described
in Section~\ref{sec:architecture}, where any network partition containing a ground block is
assigned node index~0 directly instead of a freshly allocated index. No change to the
solver itself was required: the stamping routines already handled node~0 correctly
(a term referencing node~0 is simply omitted from the assembled matrix, per the standard
MNA convention), because it was always a valid destination for a component's terminal, just
never one a player could deliberately choose. Wire and ground blocks were, in the same
change, made directly probeable, so that a player can read the absolute voltage at any wired
point once a ground reference has been established elsewhere in the same network -- turning
the abstract idea of a reference node into something a player sets up explicitly, the same
way a real measurement requires choosing where to clip the ground lead before any other
reading is meaningful.

\subsection{Adjustable sources and value propagation}

The function generator's waveform shape (sine, square, or triangle) is a preset cycled
directly on the block, but its amplitude and frequency are set externally by two small
utility blocks -- a voltage module and a frequency module -- that a player places touching
the generator, or touching a chain of other modules of either kind that eventually reaches
one or more generators. Each module tracks two independently propagated quantities, a
voltage value and a frequency value, each tagged with the game tick at which it was last
\emph{authored} (set directly by a player right-clicking that specific module). A module
only authors the channel matching its own kind; for the other channel, it purely relays
whatever value and tag it last observed. Once per tick, every module compares both of its
channels against every immediately adjacent module of either kind and adopts whichever has
the more recent authoring tick, then pushes its resolved voltage and frequency to any
adjacent generator. This lets a single control reach an arbitrarily long, possibly mixed
chain of modules and every generator at its far end, converging one hop per tick (at
20 ticks/second, imperceptibly fast for any build a player is likely to construct), while
remaining simple enough to implement without a separate, explicit ``wiring'' concept distinct
from ordinary block adjacency.
